% ============================================================
% 電圧波の観測 ワークシート本文
% 原本: 伝送線路_改v2.docx の記入欄【　】と表2・表3
% 表番号は手引書本文に合わせて 表2 から始める（表1は本文の SI 接頭語）
% ============================================================

\begin{center}
  {\sffamily\gtfamily\bfseries\Large 電気電子工学実験I　電圧波の観測　記録用紙}
\end{center}

\vspace{0.6\baselineskip}
\begin{flushright}
  学籍番号：\underline{\hspace{12\zw}}\quad
  氏名：\underline{\hspace{12\zw}}
\end{flushright}
\vspace{0.8\baselineskip}

\newcommand{\fillin}{【\hspace{6\zw}】}
\renewcommand{\arraystretch}{1.35}

\section*{実験1：電気信号の進む速さを見てみよう！}

\begin{tabular}{@{}ll@{}}
  (1) 同軸ケーブルの長さ & \fillin\ m \\
  (3) 周波数 & \fillin\ Hz \\
  (4) オシロスコープの時間軸 & \fillin\ s/Div \\
  \phantom{(4)} 測定結果：時間のずれ & \fillin\ s \\
  (5) 計算結果：電気信号伝搬速度 & \fillin\ m/s \\
  (6) 光速に対する比 $v/c$ & \fillin \\
\end{tabular}

\section*{実験2：電気信号は跳ね返る？反射波を見てみよう！}

\setcounter{table}{1}
\begin{center}
  \captionof{table}{(2)〜(4) ケーブルの端のつなぎ方による波形の違い}
  \label{tab:voltagewave-reflection}
  \begin{tabular}{|c|c|c|p{11\zw}|}
    \hline
    端のつなぎ方 & ch1 の $V_\mathrm{pp}$ (V) & ch2 の $V_\mathrm{pp}$ (V) & 波形の特徴 \\ \hline\hline
    \SI{50}{\ohm}終端 & & & \\[3mm] \hline
    開放 & & & \\[3mm] \hline
    短絡 & & & \\[3mm] \hline
  \end{tabular}
\end{center}

\begin{tabular}{@{}ll@{}}
  (5) 短絡のときの時間 & \fillin\ s \\
  (6) 開放のときの時間 & \fillin\ s \\
  (7) 実験1の時間のずれに対する比 & \fillin\ 倍 \\
\end{tabular}

\section*{実験3：定在波を見てみよう！}

\begin{tabular}{@{}ll@{}}
  (4) 各班の線の長さの合計 & \fillin\ m \\
  \phantom{(4)} その長さが波長になる周波数 & \fillin\ Hz \\
  (5) 電圧値が初めて 0.5 V を下回った位置の同軸ケーブル長 & \fillin\ m \\
\end{tabular}

\clearpage
\noindent
\begingroup
  \renewcommand{\arraystretch}{1.5}
  \begin{minipage}[t]{0.48\linewidth}
    \centering
    \captionof{table}{(5)【\hspace{4\zw}】Hz時のケーブル長に対する電圧値}
    \label{tab:voltagewave-table2}
    \begin{tabular}{|>{\centering\arraybackslash}p{7\zw}|>{\centering\arraybackslash}p{7\zw}|}
      \hline
      同軸ケーブル長 (m) & 電圧 (V) \\ \hline\hline
      & \\ \hline
      & \\ \hline
      & \\ \hline
      & \\ \hline
      & \\ \hline
      & \\ \hline
      & \\ \hline
      & \\ \hline
      & \\ \hline
      & \\ \hline
      & \\ \hline
      & \\ \hline
      & \\ \hline
      & \\ \hline
      & \\ \hline
      & \\ \hline
      & \\ \hline
      & \\ \hline
    \end{tabular}
  \end{minipage}\hfill
  \begin{minipage}[t]{0.48\linewidth}
    \centering
    \captionof{table}{(6)【\hspace{4\zw}】Hz時のケーブル長に対する電圧値}
    \label{tab:voltagewave-table3}
    \begin{tabular}{|>{\centering\arraybackslash}p{7\zw}|>{\centering\arraybackslash}p{7\zw}|}
      \hline
      同軸ケーブル長 (m) & 電圧 (V) \\ \hline\hline
      & \\ \hline
      & \\ \hline
      & \\ \hline
      & \\ \hline
      & \\ \hline
      & \\ \hline
      & \\ \hline
      & \\ \hline
      & \\ \hline
      & \\ \hline
      & \\ \hline
      & \\ \hline
      & \\ \hline
      & \\ \hline
      & \\ \hline
      & \\ \hline
      & \\ \hline
      & \\ \hline
    \end{tabular}
  \end{minipage}
\endgroup
